\documentclass[11pt]{article}
\usepackage{fontspec,amsmath,amssymb,amsthm,hyperref}
\usepackage[margin=1in]{geometry}
\setlength{\emergencystretch}{3em}
\title{Exact Gibbs invariance of the actual Hamiltonian flow}
\author{LangevinLean mathematical exposition}
\date{}
\begin{document}\maketitle
\noindent Source module: \texttt{KineticLangevin/HamiltonianStationarity.lean}. Compile this proof document with XeLaTeX or LuaLaTeX. The accompanying HTML pairs the mathematical argument with the exact Lean statements and records the scope of the formalization.


This module closes the deterministic Hamiltonian stationarity argument. Explicit corrected Verlet proposals preserve the Gibbs measure exactly. Their rejection probability is quadratically small, while the uncorrected Verlet iterates converge to the actual selected Hamiltonian integral solution. Bounded tests transfer invariance through this limit. The finite-dimensional strong-convexity theorem derives every required tail premise, so its statement assumes no invariant flow, numerical convergence, or error contract.


\section{Make the actual endpoint a measurable function of its initial state}


The constructed kinetic endpoint is jointly continuous in initialization and continuous forcing. Fix zero forcing, friction zero, inverse mass one, and endpoint time \(T\). Composing this joint map with \(z\mapsto(z,0)\) proves continuity of the selected Hamiltonian endpoint. Lean supplies the compact forcing-space zero explicitly and simplifies the composed definition. The endpoint pushforward is therefore a concrete measurable law.



\noindent Lean declarations: \texttt{hamiltonianFlow\_endpoint\_continuous}.

\section{Pass exact corrected invariance to the actual integral solution}


Let \(\pi\) be the actual Hamiltonian weighted measure, \(B=\int(\|v\|^2+\|\nabla f(x)\|^2)\,d\pi\), and \(K=M(2M+11)/8\). Use \(N=n+1\) Verlet steps of size \(h=T/N\). Exact corrected invariance and the iterated bounded-test estimate imply, for \(\|F\|\le1\), \[\left\|\int F(V_h^Nz)\,d\pi-\int F\,d\pi\right\|\le 2N\int(1-a_h)\,d\pi\le\frac{2T^2KB}{N}.\] The mesh is eventually at most one, so the proved quadratic rejection estimate applies; the displayed bound tends to zero. The independently proved numerical theorem gives \(V_{T/N}^Nz\to\Phi_Tz\), where \(\Phi\) is the actual Picard integral solution. The bounded-test limit theorem identifies \((\Phi_T)_\#\pi=\pi\). Lean assembles these actual results, proves the mesh limit, checks nonnegative multipliers, and simplifies the exact scalar bound. It assumes only the finite measure and precisely displayed quadratic moment required by this intermediate general theorem.



\noindent Lean declarations: \texttt{hamiltonianFlow\_preserves\_gibbs\_of\_moment}.

\section{Discharge every tail premise from the potential assumptions}


In finite dimensions, positive strong convexity and differentiability make the actual Hamiltonian Gibbs measure finite by the factorization and Gaussian domination proved in the moments module. The globally Lipschitz gradient supplies the finite force second moment, and the thermal velocity factor supplies the velocity second moment. Substituting these checked results into the preceding limit theorem gives exact Hamiltonian Gibbs invariance from the original first-order potential assumptions. The theorem applies to every nonnegative finite horizon. It proves the Hamiltonian subflow; the full friction-and-noise kinetic diffusion additionally needs its thermostat splitting limit.



\noindent Lean declarations: \texttt{hamiltonianFlow\_preserves\_gibbs}.

\end{document}
